Um die Beziehung \( S^{-1}(\Lambda) = \left(\gamma^0 S(\Lambda) \gamma^0\right)^{\dagger} \) zu zeigen, beginnen wir mit der Definition von \( S(\Lambda) \) und den Eigenschaften der Dirac-Matrizen. Wir ergänzen die Herleitung der Beziehung für antihermites \( A \).

\begin{align*}
S(\Lambda) &= e^{-\frac{i}{4} \omega^{\mu \nu} \sigma_{\mu \nu}}, \\
\sigma_{\mu \nu} &= \frac{i}{2} [\gamma_\mu, \gamma_\nu].
\end{align*}

Die Inverse von \( S(\Lambda) \) ist gegeben durch:

\begin{align*}
S^{-1}(\Lambda) &= e^{\frac{i}{4} \omega^{\mu \nu} \sigma_{\mu \nu}}.
\end{align*}

Für eine antihermites Matrix \( A \) gilt \( A^{\dagger} = -A \). Daraus folgt:

\begin{align*}
(e^A)^{\dagger} &= \left(\sum_{n=0}^{\infty} \frac{A^n}{n!}\right)^{\dagger} = \sum_{n=0}^{\infty} \frac{(A^n)^{\dagger}}{n!} = \sum_{n=0}^{\infty} \frac{(A^{\dagger})^n}{n!} \\
&= \sum_{n=0}^{\infty} \frac{(-A)^n}{n!} = e^{-A} = e^{A^{\dagger}}.
\end{align*}

Da \( \sigma_{\mu \nu} \) antihermitisch ist, folgt:

\begin{align*}
(\gamma^0 \sigma_{\mu \nu} \gamma^0) &= \sigma_{\mu \nu}^\dagger.
\end{align*}

Daraus ergibt sich:

\begin{align*}
\left(\gamma^0 S(\Lambda) \gamma^0\right) &= e^{-\frac{i}{4} \omega^{\mu \nu} (\gamma^0 \sigma_{\mu \nu} \gamma^0)} = e^{-\frac{i}{4} \omega^{\mu \nu} \sigma_{\mu \nu}^\dagger}.
\end{align*}

Anwendung der Exponentialregel für antihermites \( A \):

\begin{align*}
\left(\gamma^0 S(\Lambda) \gamma^0\right)^{\dagger} &= \left(e^{-\frac{i}{4} \omega^{\mu \nu} \sigma_{\mu \nu}^\dagger}\right)^{\dagger} = e^{\frac{i}{4} \omega^{\mu \nu} \sigma_{\mu \nu}}.
\end{align*}

Dies zeigt, dass:

\begin{align*}
S^{-1}(\Lambda) &= \left(\gamma^0 S(\Lambda) \gamma^0\right)^{\dagger}.
\end{align*}

%CHECK: M1 wurde durch die explizite Herleitung der Exponentialregel für antihermites A adressiert.